% GENERATED FILE. Edit canonical lessons, then run npm run generate:books.
% canonical-source-hash: fnv1a64:c7acbf88f855d9e7
% canonical-lessons: TA-C377-vanmurai, TA-C377-akimcai, TA-C377-orrumai, TA-C377-verrumai, TA-C377-pilavupatu

\chapter{Unity and Division}
\label{ch:ta-unity-and-division}
\hlchaptermodality{377}

\begin{chapteropening}
\textbf{By the end of this chapter:} \emph{I can talk about violence, non-violence and unity.}

Complete the last lesson of chapter 377: I can talk about violence, non-violence and unity.
\end{chapteropening}

\section[\texorpdfstring{\ta{வன்முறை} (vaṉmuṟai)}{வன்முறை (vaṉmuṟai)}]{\ta{வன்முறை} (vaṉmuṟai) — violence}

\label{lesson:TA-C377-vanmurai}



\begin{quote}
Before the new one: say the Tamil for to demand, then the Tamil for a petition.
\end{quote}

\subsection*{You'll want to know: \ta{வன்முறை}}
\textbf{\ta{வன்முறை}} — \emph{vaṉmuṟai} — \textquotedblleft{}violence\textquotedblright{}.

\begin{quote}
\ta{நாங்கள்} \ta{வன்முறையை} \ta{ஏற்றுக்கொள்ள} \ta{மாட்டோம்}.
\end{quote}

\begin{quote}
\emph{nāṅkaḷ vaṉmuṟaiyai ēṟṟukkoḷḷa māṭṭōm.}
\end{quote}

\begin{quote}
We will not accept violence.
\end{quote}

\ta{வல்}, hard + \ta{முறை}, way: the hard way.

\begin{cousinweb}
\textbf{\ta{வன்முறை}} (\emph{vaṉmuṟai}) is native: \emph{val}, \textquotedblleft{}strong, hard, forceful\textquotedblright{} (DEDR 5276), whose l becomes \ta{ன்} before m, and \emph{muṟai}, \textquotedblleft{}order, way\textquotedblright{} (DEDR 5015).
\end{cousinweb}

\begin{grammarlens}[title={What you've built}]
One more word for giving reasons and opinions.
\end{grammarlens}

\subsection*{Your turn}
\begin{itemize}
\raggedright
  \item \emph{Say it:} \emph{vaṉmuṟai}
  \item \emph{Say it:} \emph{vaṉmuṟai}, once more
  \item \emph{Say it:} \emph{nāṅkaḷ vaṉmuṟaiyai ēṟṟukkoḷḷa māṭṭōm.} — We will not accept violence.
  \item \emph{Recall:} say the Tamil for to demand, then the Tamil for a petition, then say \emph{vaṉmuṟai} again
\end{itemize}

\begin{tcolorbox}[breakable,colback=teal!4,colframe=teal!35!black,title={Before you move on}]
What does \ta{வன்முறை} mean? (\textquotedblleft{}Violence\textquotedblright{}.) Say it once more.
\end{tcolorbox}
\section[\texorpdfstring{\ta{அகிம்சை} (akimcai)}{அகிம்சை (akimcai)}]{\ta{அகிம்சை} (akimcai) — non-violence}

\label{lesson:TA-C377-akimcai}



\begin{quote}
Before the new one: say the Tamil for a petition, then the Tamil for violence.
\end{quote}

\subsection*{You'll want to know: \ta{அகிம்சை}}
\textbf{\ta{அகிம்சை}} — \emph{akimcai} — \textquotedblleft{}non-violence\textquotedblright{}.

\begin{quote}
\ta{அவர்} \ta{அகிம்சை} \ta{வழியில்} \ta{போராடினார்}.
\end{quote}

\begin{quote}
\emph{avar akimcai vaḻiyil pōrāṭiṉār.}
\end{quote}

\begin{quote}
He fought in the way of non-violence.
\end{quote}

\ta{அகிம்சை} \ta{வழி}, the path of non-violence.

\begin{cousinweb}
\textbf{\ta{அகிம்சை}} (\emph{akimcai}) is borrowed from Sanskrit \emph{ahiṃsā}, \textquotedblleft{}not injuring anything\textquotedblright{} (Monier-Williams). Tamil writes the Sanskrit h as \ta{க} and the final -ā as -\ta{ஐ}.
\end{cousinweb}

\begin{grammarlens}[title={What you've built}]
One more word for giving reasons and opinions.
\end{grammarlens}

\subsection*{Your turn}
\begin{itemize}
\raggedright
  \item \emph{Say it:} \emph{akimcai}
  \item \emph{Say it:} \emph{akimcai}, once more
  \item \emph{Say it:} \emph{avar akimcai vaḻiyil pōrāṭiṉār.} — He fought in the way of non-violence.
  \item \emph{Recall:} say the Tamil for a petition, then the Tamil for violence, then say \emph{akimcai} again
\end{itemize}

\begin{tcolorbox}[breakable,colback=teal!4,colframe=teal!35!black,title={Before you move on}]
What does \ta{அகிம்சை} mean? (\textquotedblleft{}Non-violence\textquotedblright{}.) Say it once more.
\end{tcolorbox}
\section[\texorpdfstring{\ta{ஒற்றுமை} (oṟṟumai)}{ஒற்றுமை (oṟṟumai)}]{\ta{ஒற்றுமை} (oṟṟumai) — unity}

\label{lesson:TA-C377-orrumai}



\begin{quote}
Before the new one: say the Tamil for violence, then the Tamil for non-violence.
\end{quote}

\subsection*{You'll want to know: \ta{ஒற்றுமை}}
\textbf{\ta{ஒற்றுமை}} — \emph{oṟṟumai} — \textquotedblleft{}unity\textquotedblright{}.

\begin{quote}
\ta{எங்கள்} \ta{ஊர்} \ta{மக்களிடையே} \ta{நல்ல} \ta{ஒற்றுமை} \ta{இருக்கிறது}.
\end{quote}

\begin{quote}
\emph{eṅkaḷ ūr makkaḷiṭaiyē nalla oṟṟumai irukkiṟatu.}
\end{quote}

\begin{quote}
There is great unity among the people of our town.
\end{quote}

\ta{ஒற்றுமையாக}, united: \ta{ஒற்றுமையாக} \ta{வாழ்}, to live in harmony.

\begin{cousinweb}
\textbf{\ta{ஒற்றுமை}} (\emph{oṟṟumai}) is native Dravidian: DEDR 990 lists Tamil \emph{oru}, \textquotedblleft{}one\textquotedblright{}, and \emph{oṟṟumai}, \textquotedblleft{}union, oneness\textquotedblright{}.
\end{cousinweb}

\begin{grammarlens}[title={What you've built}]
One more word for giving reasons and opinions.
\end{grammarlens}

\subsection*{Your turn}
\begin{itemize}
\raggedright
  \item \emph{Say it:} \emph{oṟṟumai}
  \item \emph{Say it:} \emph{oṟṟumai}, once more
  \item \emph{Say it:} \emph{eṅkaḷ ūr makkaḷiṭaiyē nalla oṟṟumai irukkiṟatu.} — There is great unity among the people of our town.
  \item \emph{Recall:} say the Tamil for violence, then the Tamil for non-violence, then say \emph{oṟṟumai} again
\end{itemize}

\begin{tcolorbox}[breakable,colback=teal!4,colframe=teal!35!black,title={Before you move on}]
What does \ta{ஒற்றுமை} mean? (\textquotedblleft{}Unity\textquotedblright{}.) Say it once more.
\end{tcolorbox}
\section[\texorpdfstring{\ta{வேற்றுமை} (vēṟṟumai)}{வேற்றுமை (vēṟṟumai)}]{\ta{வேற்றுமை} (vēṟṟumai) — diversity; difference}

\label{lesson:TA-C377-verrumai}



\begin{quote}
Before the new one: say the Tamil for non-violence, then the Tamil for unity.
\end{quote}

\subsection*{You'll want to know: \ta{வேற்றுமை}}
\textbf{\ta{வேற்றுமை}} — \emph{vēṟṟumai} — \textquotedblleft{}diversity; difference\textquotedblright{}.

\begin{quote}
\ta{வேற்றுமையில்} \ta{ஒற்றுமை} \ta{நம்} \ta{நாட்டின்} \ta{சிறப்பு}.
\end{quote}

\begin{quote}
\emph{vēṟṟumaiyil oṟṟumai nam nāṭṭiṉ ciṟappu.}
\end{quote}

\begin{quote}
Unity in diversity is what makes our country special.
\end{quote}

In grammar, \ta{வேற்றுமை} also names the cases of a noun.

\begin{cousinweb}
\textbf{\ta{வேற்றுமை}} (\emph{vēṟṟumai}) is native Dravidian: DEDR 5548 lists Tamil \emph{vēṟu}, \textquotedblleft{}other, different\textquotedblright{}, and \emph{vēṟṟumai}, \textquotedblleft{}difference\textquotedblright{}.
\end{cousinweb}

\begin{grammarlens}[title={What you've built}]
One more word for giving reasons and opinions.
\end{grammarlens}

\subsection*{Your turn}
\begin{itemize}
\raggedright
  \item \emph{Say it:} \emph{vēṟṟumai}
  \item \emph{Say it:} \emph{vēṟṟumai}, once more
  \item \emph{Say it:} \emph{vēṟṟumaiyil oṟṟumai nam nāṭṭiṉ ciṟappu.} — Unity in diversity is what makes our country special.
  \item \emph{Recall:} say the Tamil for non-violence, then the Tamil for unity, then say \emph{vēṟṟumai} again
\end{itemize}

\begin{tcolorbox}[breakable,colback=teal!4,colframe=teal!35!black,title={Before you move on}]
What does \ta{வேற்றுமை} mean? (\textquotedblleft{}Diversity; difference\textquotedblright{}.) Say it once more.
\end{tcolorbox}
\section[\texorpdfstring{\ta{பிளவுபடு} (piḷavupaṭu)}{பிளவுபடு (piḷavupaṭu)}]{\ta{பிளவுபடு} (piḷavupaṭu) — to be divided; to split}

\label{lesson:TA-C377-pilavupatu}



\begin{quote}
Before the new one: say the Tamil for unity, then the Tamil for diversity.
\end{quote}

\subsection*{You'll want to know: \ta{பிளவுபடு}}
\textbf{\ta{பிளவுபடு}} — \emph{piḷavupaṭu} — \textquotedblleft{}to be divided; to split\textquotedblright{}.

\begin{quote}
\ta{அந்த} \ta{முடிவால்} \ta{எங்கள்} \ta{குழு} \ta{இரண்டாகப்} \ta{பிளவுபட்டது}.
\end{quote}

\begin{quote}
\emph{anta muṭivāl eṅkaḷ kuḻu iraṇṭākap piḷavupaṭṭatu.}
\end{quote}

\begin{quote}
Because of that decision, our group split in two.
\end{quote}

Past (it …): \ta{பிளவுபட்டது}, \emph{piḷavupaṭṭatu}.

\ta{பிளவு}, a split + \ta{படு}, to undergo.

\begin{cousinweb}
\textbf{\ta{பிளவுபடு}} (\emph{piḷavupaṭu}) is native: DEDR 4194 lists Tamil \emph{piḷa}, \textquotedblleft{}to be split\textquotedblright{}, and \emph{piḷavu}, \textquotedblleft{}cleft, division, disunion\textquotedblright{}, here with \emph{paṭu}, \textquotedblleft{}to undergo\textquotedblright{}.
\end{cousinweb}

\begin{grammarlens}[title={What you've built}]
One more word for giving reasons and opinions.
\end{grammarlens}

\subsection*{Your turn}
\begin{itemize}
\raggedright
  \item \emph{Say it:} \emph{piḷavupaṭu}
  \item \emph{Say it:} \emph{piḷavupaṭu}, once more
  \item \emph{Say it:} \emph{anta muṭivāl eṅkaḷ kuḻu iraṇṭākap piḷavupaṭṭatu.} — Because of that decision, our group split in two.
  \item \emph{Recall:} say the Tamil for unity, then the Tamil for diversity, then say \emph{piḷavupaṭu} again
\end{itemize}

\begin{tcolorbox}[breakable,colback=teal!4,colframe=teal!35!black,title={Before you move on}]
What does \ta{பிளவுபடு} mean? (\textquotedblleft{}To be divided; to split\textquotedblright{}.) Say it once more.
\end{tcolorbox}
